\documentclass{article}
\usepackage{logicproof}
\usepackage{amssymb}
\usepackage{amsmath}

\begin{document}

\section*{Question 2: Natural Deduction Proofs}

\subsection*{a) $\forall x (P(x) \rightarrow Q(x)), \exists x P(x) \vdash \exists x Q(x)$}

\begin{logicproof}{2}
    \forall x (P(x) \rightarrow Q(x)) & premise \\
    \exists x P(x) & premise \\
    \begin{subproof}
        x_0 \quad P(x_0) & assume \\
        P(x_0) \rightarrow Q(x_0) & $\forall e$ 1 \\
        Q(x_0) & $\rightarrow e$ 4, 3 \\
        \exists x Q(x) & $\exists i$ 5
    \end{subproof}
    \exists x Q(x) & $\exists e$ 2, 3--6
\end{logicproof}

\subsection*{b) $\forall x P(x) \lor \forall x Q(x) \vdash \forall x (P(x) \vee Q(x))$}

\begin{logicproof}{3} %3 is the number of assumptions
    \forall x P(x) \lor \forall x Q(x) & premise \\
    \begin{subproof}
        \forall x P(x) & assume \\
        \begin{subproof}
            x_0 & arbitrary \\
            P(x_0) & $\forall e$ 2 \\
            P(x_0) \lor Q(x_0) & $\lor i_1$ 4
        \end{subproof}
        \forall x (P(x) \lor Q(x)) & $\forall i$ 3--5
    \end{subproof}
    \begin{subproof}
        \forall x Q(x) & assume \\
        \begin{subproof}
            y_0 & arbitrary \\
            Q(y_0) & $\forall e$ 7 \\
            P(y_0) \lor Q(y_0) & $\lor i_2$ 9
        \end{subproof}
        \forall x (P(x) \lor Q(x)) & $\forall i$ 8--10
    \end{subproof}
    \forall x (P(x) \lor Q(x)) & $\lor e$ 1, 2--6, 7--11
\end{logicproof}

\newpage
\section*{Question 3: Model Finding}
\subsection*{a) $\forall x \exists y P(f(x), y)$}
\textbf{Model that satisfies ($\mathcal{M} \models \phi$):}
\begin{itemize}
    \item Domain $D = \mathbb{Z}$ (integers).
    \item $f(x) = x$ (identity function).
    \item $P(x, y)$ is defined as the relation $x < y$.
    \item \textbf{Explanation:} For every integer $x$, there exists an integer $y$ (such as $y = x+1$) such that $x < y$ holds.
\end{itemize}

\textbf{Model that does not satisfy ($\mathcal{M} \not\models \phi$):}
\begin{itemize}
    \item Domain $D = \{1, 2\}$.
    \item $f(x) = 1$ for all $x$.
    \item $P(x, y)$ is always \textit{false}.
    \item \textbf{Explanation:} Since $P(1, 1)$ and $P(1, 2)$ are both false, there is no $y$ that satisfies the relation for any $x$.
\end{itemize}

\subsection*{b) $(\forall x P(x) \rightarrow S) \vdash \exists x (P(x) \rightarrow S)$}
\textbf{Model that satisfies:}
\begin{itemize}
    \item Domain $D = \{1\}$.
    \item $P(1) = \text{true}$.
    \item $S = \text{true}$.
    \item \textbf{Explanation:} The premise $(\forall x P(x) \rightarrow S)$ is $\text{true} \rightarrow \text{true}$, which is true. The conclusion $\exists x (P(x) \rightarrow S)$ is true because $P(1) \rightarrow S$ is true.
\end{itemize}

\textbf{Model that does not satisfy (Counterexample):}
\begin{itemize}
    \item Domain $D = \emptyset$ (The Empty Set).
    \item \textbf{Explanation:} In standard first-order logic, the domain must be non-empty. However, if we consider a non-empty domain $D=\{1\}$ where:
    \item $P(1) = \text{false}$ and $S = \text{false}$.
    \item The premise $(\forall x P(x) \rightarrow S)$ becomes $(\text{false} \rightarrow \text{false})$, which is \textbf{true}.
    \item The conclusion $\exists x (P(x) \rightarrow S)$ becomes $(\text{false} \rightarrow \text{false})$, which is also \textbf{true}.
    \item \textit{Note:} This specific sequent is actually a valid logical principle in classical logic with non-empty domains, meaning a counterexample would only exist if the domain were empty, which is usually disallowed.
\end{itemize}

\end{document}